\chapter{Calibration and threshold derivation}

The phase where most of the work turned out to be arranging for numbers to be
checkable rather than producing them.

\begin{defect}{A threshold definition whose answer is always zero}
\symptom{The low threshold was specified as the smallest confidence at which the
near miss band still captures at least half of the recall the high threshold
gives up. Written down carefully, that has a degenerate answer.}
\diagnosis{Captured recall only ever falls as the threshold rises, so every value
below a qualifying one also qualifies, the smallest qualifying value is always
zero, and a threshold of zero says nothing.}
\resolution{The non degenerate reading is the \emph{greatest} qualifying value,
and it was written into the prediction file before the model existed.}
\instead{Picking whichever reading produced a nicer boundary once both were
computable. Writing the reading down first settled a real ambiguity at a point
where there were not yet two candidate answers to choose between, which is the
only time that choice is free.}
\end{defect}

\begin{defect}{One threshold document would have moved every rule engine finding}
\symptom{The derived n-gram boundary sits far above the rule engine's confident
band. A single pair of thresholds shared by both engines would have moved every
rule engine finding on every page, churned every golden snapshot, and done it as
a side effect of training a model.}
\diagnosis{A confidence scale is a property of an engine. A calibrated
probability and a rule tier band are different quantities that happen to be
written in the same units.}
\resolution{The threshold document carries a block per engine at a bumped schema
version, the loader takes an engine name, and the command line resolves the
engine first and then asks for that engine's block. The rule block is byte
identical to what the previous phase committed, and the golden diff after the
whole phase is the version string and nothing else, which is the evidence that it
worked.}
\instead{One global pair, or an inheritance rule for an engine with no block. An
engine with no block raises. That was designed to bite when the language model
arrived, and it did, which is how that engine's threshold decision came to be
made deliberately and in advance rather than by inheriting a boundary derived on
a different scale.}
\end{defect}

\begin{defect}{A precision of one over a denominator nobody could see}
\symptom{The derived block reported that the precision target was met.}
\diagnosis{A precision of one over fifteen accusations and a precision of one
over fifteen hundred are the same number and are not the same claim. A reader of
the file could not tell which they were looking at.}
\resolution{The derived block records the denominators beside the rates, and the
model card says the same thing in words.}
\instead{Reporting the rate alone and mentioning the denominator in prose
somewhere else. The file is what the runtime reads and what a reader opens
first.}
\end{defect}

\begin{defect}{Calibration made the summary statistic worse}
\symptom{Expected calibration error was lower before calibration than after it
(\vEceBefore{} against \vEceAfter{}).}
\diagnosis{The uncalibrated model on this corpus is already close to honest.
Refitting a one versus rest sigmoid per class on a development split of one
template per family, and then renormalising, moves classes that were already well
behaved.}
\resolution{The calibrator stays, and the model card says the error went the
wrong way and by how much. Two reasons: macro-F1 is higher with it
(\vMacroFBeforeCalibration{} before against \vDevMacroF{} after), and the
argument for calibrating at all is about the number a single finding quotes to a
developer rather than about an average gap across a split.}
\instead{Dropping the calibrator. Removing a component because its summary
statistic got worse, without re-registering a prediction first, is tuning on the
thing being measured. The honest move was to keep the design, publish the
number that went the wrong way, and say why the design still stands.}
\end{defect}

\begin{defect}{The sweep chose the edge of its own grid}
\symptom{The chosen inverse regularisation is the largest value the
pre-registered grid offered (Table in the main report's classifier chapter).}
\diagnosis{Either the grid is too narrow or the model wants less regularisation
than expected, and the data cannot distinguish those without a wider grid.}
\resolution{The grid was not widened. The consequence is recorded in the model
card instead: this model may be less regularised than a wider grid would have
chosen.}
\instead{Widening the grid and re-running. Widening a pre-registered grid because
the answer landed at its edge is how a sweep becomes a search for a number. A
later phase that revisits it re-registers first.}
\end{defect}

\begin{defect}{An abstention that was right by accident}
\symptom{A model emitting the unknown label with a high calibrated probability
reached the correct outcome, but through the wrong mechanism: the declaration
lookup returned nothing, rather than the confidence gating anything.}
\diagnosis{Right answer, wrong reason, and a wrong reason is a defect waiting for
the code around it to change.}
\resolution{An explicit branch. When the winning class is the unknown label the
reported confidence is zero, exactly as the rule engine reports it, and the
calibrated probability of the runner up stays on its own field where the
confusion analysis can still read it.}
\instead{Leaving it, on the grounds that the outcome was correct. The outcome was
correct for a reason that no test asserted and that any refactor of the
declaration lookup would have silently removed.}
\end{defect}

\begin{defect}{The model's first wrong accusations}
\symptom{After the corpus repair, the n-gram engine accused
\vNgramMissingAccusations{} fields of a missing declaration and was wrong about
\vNgramMissingWrong{} of them. That is the first time in this project that an
engine told a developer to add a token the answer key disagrees with.}
\diagnosis{Not a regression in the threshold policy. The target precision behind
the high threshold is the pre-registered one and the derivation script was run
unmodified. What changed is that the derived threshold now sits on a model whose
development precision at that point was itself below one, so a test split
precision below one is the design working rather than failing.}
\resolution{Recorded rather than rounded, in the engineering log, the model card
and the main report. A precision target is a target, not a guarantee.}
\instead{Raising the target until the wrong accusations disappeared. That is
tuning a pre-registered constant against the test split, which is the single
worst thing available in this project's rulebook.}
\end{defect}

\begin{defect}{A class that finally reached the isotonic switchover}
\symptom{Not a defect. A prediction that was correct for one corpus and stopped
being correct for the next.}
\diagnosis{The earlier phase predicted that no class would ever have enough
development positives for isotonic regression, and for its corpus that was true
of every class. After the repair, the label covering search boxes and consent
checkboxes crossed the switchover.}
\resolution{The per class method is recorded in the model card, so the change is
visible rather than implicit.}
\instead{Nothing to refuse here. It is in this chapter because a prediction that
holds under one corpus and fails under a larger one is exactly the kind of thing
that gets quietly forgotten, and the record of it is what makes the next such
prediction more careful.}
\end{defect}
